% !TEX program = xelatex
\documentclass[11pt,a4paper]{ctexart}

% ── 页面 ──
\usepackage[margin=2.2cm]{geometry}
\usepackage{hyperref}
\usepackage{enumitem}
\usepackage{tabularx}
\usepackage{booktabs}
\usepackage{graphicx}
\usepackage{float}

\hypersetup{
    colorlinks=true, linkcolor=blue, urlcolor=blue,
    pdftitle={SJTU Agent — 来自学生开发者的合作邀请},
    pdfauthor={SJTU Agent 开发团队},
}
\setlist{nosep,leftmargin=*,topsep=3pt,itemsep=1pt}

\begin{document}

\begin{center}
    {\LARGE \textbf{SJTU Agent}}
    \vspace{0.3em}

    {\large 来自学生开发者的合作邀请}
    \vspace{0.5em}

    \small 2026年6月 \quad·\quad v0.6.3 \quad·\quad \href{https://github.com/kuan-er/sjtu-agent}{github.com/kuan-er/sjtu-agent}
\end{center}

\vspace{0.5em}
\hrule
\vspace{0.5em}

\begin{quote}\small
2026 年五一假期，几个交大学生开始写一个校园 AI 助手。45 天、12 个版本、69 个 Star 之后，我们想和学校聊聊。
\end{quote}

\section*{项目速览}

\begin{tabularx}{\textwidth}{lX}
\textbf{启动时间} & 2026 年 5 月 5 日 \\
\textbf{当前版本} & v0.6.2（12 次发布）\\
\textbf{代码规模} & \textasciitilde20,000 行 Python \\
\textbf{测试覆盖} & 287 个自动化测试，GitHub Actions CI \\
\textbf{贡献者} & Azzygoatcoder（300 提交）、kuan-er（103 提交）、dajiaohuang、Aixiaoxia、JackyTJie \\
\textbf{许可证} & MIT 开源 \\
\end{tabularx}

\vspace{0.3em}
\textbf{学生维护、开源、零经费。项目主要通过 Claude Code 等 Coding Agent 辅助开发——AI 时代，学生参与校园信息化的门槛已经大幅降低。}

\section{安全性（合作的前提）}

以下措施已在代码中实施并持续维护：
\begin{itemize}
    \item 用户凭据储存在本地，\textbf{不上传任何远程服务器}
    \item 凭据文件自动设 \texttt{0o600} 权限，配置文件原子写入防崩溃
    \item 飞书 Bot 白名单控制，所有命令均需通过校验
    \item Python 执行时自动剥离 \texttt{JACCOUNT\_PASSWORD} 等敏感环境变量
    \item 外部请求强制 HTTPS，URL 校验防 SSRF（私有 IP、回环地址）
    \item 文件路径 \texttt{resolve() + is\_relative\_to()} 防穿越
\end{itemize}

\section{为什么需要官方支持}

我们\textbf{打通了所有能通过公开渠道获取的数据源}。但以下能力，学生力量做不到：

\vspace{0.3em}
\begin{tabularx}{\textwidth}{lX}
\toprule
\textbf{重要度} & \textbf{想做但做不了的事} \\
\midrule
{[重要]} & 食堂推荐接入一码通消费记录（脱敏） \\
{[重要]} & 校历自动获取，替代手动录入 JSON \\
{[重要]} & 稳定的 jAccount OAuth，替代 Playwright 模拟登录 \\
{[中等]} & 课表/成绩的教务系统数据接口 \\
{[中等]} & 学校云平台托管，支持模块化安装 \\
{[一般]} & 移动端入口——接入“交我办”小程序 \\
\bottomrule
\end{tabularx}

\subsection*{一个例子：\texttt{/eat}}

campuslife API 是公开接口，所以食堂推荐做出来了。但如果能接入一码通消费记录（脱敏），推荐精度能再上一个量级。贡献者 @JackyTJie 尝试抓包后发现：
\begin{quote}\small
“全都加密了，认证非常严格，仅限于交我办软件访问。除非先获得官方支持，不然做不出来。”
\end{quote}

\subsection*{大模型怎么办}

三种方案：\textbf{（1）致远一号}——交大已有免费 LLM 平台；\textbf{（2）学生自选 API Key}——费用透明；\textbf{（3）学校统一采购}——全校基础额度。倾向方案 1 或 2，不对学校产生额外经费压力。

\subsection*{模块化安装}

接入学校平台后，同学选择需要的功能，系统按需拉起服务，不强制全装。

\section{已经做了什么}

\begin{tabularx}{\textwidth}{lX}
\toprule
\textbf{数据源} & \textbf{获取方式} \\
\midrule
Canvas LMS & API（需用户手动获取 Token）\\
教学信息服务网 & 模拟登录 + 爬虫 \\
物理实验平台 phycai & 模拟登录 + 爬虫 \\
中国大学 MOOC / AI 好课 & 模拟登录 + 爬虫 \\
campuslife API & \textbf{公开接口（无需认证）}\\
教务处通知、水源社区 & 网页抓取 \\
交大邮箱 & IMAP（readonly）\\
\bottomrule
\end{tabularx}

\subsection*{飞书 Bot 主要命令}

\begin{verbatim}
/eat     实时拥挤度 + 个性化食堂推荐 + 偏好学习
/news    校园新闻 AI 摘要
/hw do 3 下载并分析第 3 个作业（Canvas → LLM → 解答）
/template 一键套用 SJTU 毕业论文 LaTeX 模板
/aihot   今日 AI 圈精选新闻
\end{verbatim}

\begin{figure}[H]
\centering
\includegraphics[width=0.48\textwidth]{docs/images/feishu-eat.png}\hfill
\includegraphics[width=0.48\textwidth]{docs/images/feishu-news.png}
\caption{\small \texttt{/eat} 食堂推荐（左）与 \texttt{/news} 校园新闻（右）}
\end{figure}

\begin{figure}[H]
\centering
\includegraphics[width=\textwidth]{docs/images/arch-diagram.png}
\caption{\small 系统架构：校园数据源 → SJTU Agent → 飞书/微信/QQ}
\end{figure}

\section{对学校的价值}

\begin{enumerate}
    \item \textbf{零研发成本}——已完成开发，学校无需投入
    \item \textbf{需求真实}——食堂推荐（午饭排队）、DDL 聚合（忘截止日期）、校历识别（节假日推送）
    \item \textbf{技术可控}——纯 Python，MIT 开源，数据本地存储，287 测试
    \item \textbf{已有基础}——69 Stars，11 Forks，5 位贡献者，45 天 20+ PR
\end{enumerate}

\section{我们的底线与下一步}

\textbf{帮同学省时间，不做越界的事。}不做考勤监控、不分析行为、不推送广告。Bot 可以提醒“今天有 3 个 DDL 截止”，但怎么做是同学自己的事。所有功能选项掌握在学生手里——屏蔽某类新闻（\texttt{/news\_block}）、重置推荐画像、关闭日报推送。

\textbf{数据主权在学生。}当前所有对话记录、用餐偏好、课表缓存存储在本地。如果未来接入学校平台，我们同样坚持：学生可以随时导出、随时删除、可以选择不参与任何数据收集。

我们不是来要经费的。只是想约一个时间，和信息化办公室的老师聊聊——展示一下我们在做什么，听听老师的建议，探讨有没有合作的可能。

\vspace{0.5em}
\begin{center}
\href{https://github.com/kuan-er/sjtu-agent}{github.com/kuan-er/sjtu-agent}
\end{center}

\end{document}
